\documentclass{article}
\usepackage[utf8]{inputenc}
\usepackage[T1]{fontenc}
\usepackage{geometry}
\geometry{
    a4paper,
    left=2.5cm,
    right=2.5cm,
    top=2.5cm,
    bottom=2.5cm
}
\usepackage{graphicx}
\usepackage{float}
\usepackage{hyperref}
\usepackage[english]{babel}
\usepackage{amsmath}
\usepackage{amssymb}
\usepackage{array}
\usepackage{booktabs}
\usepackage{caption}
\usepackage{subcaption}
\usepackage{enumitem}
\usepackage{fancyhdr}
\usepackage{titlesec}
\usepackage{url}
\usepackage{textcomp}
\usepackage{xcolor}
\usepackage{listings}
\usepackage{multirow}
\usepackage{longtable}
\usepackage{tabularx}
\usepackage{times}

% Set up page style
\pagestyle{fancy}
\fancyhf{}
\rfoot{\thepage}
\renewcommand{\headrulewidth}{0pt}

% Title formatting
\titleformat{\chapter}[display]{\normalfont\bfseries}{\chaptertitle\thechapter}{0pt}{\Huge}
\titlespacing*{\chapter}{0pt}{50pt}{40pt}

% Document info
\title{Template Document}
\author{Academic Support Agent}
\date{\today}

\begin{document}
\maketitle

\maketitle

Once upon a time\ldots This document shows how you can get ePub-like formatting in \LaTeX{} with the \verb|memoir| document class. You can't yet export directly to ePub from writeLaTeX, but you can download the source and run it through a format conversion tool, such as \verb|htlatex| to get HTML, and then go from HTML to ePub with a tool like Sigil or Calibre. See \url{http://tex.stackexchange.com/questions/16569} for more advice. And they lived happily ever after.

\lipsum
\end{document}